% ==========================================================================
% Supervisor Meeting Agenda - 16 March 2026 (30 minutes)
%
% Compile with: pdflatex agenda.tex
% ==========================================================================

\documentclass[11pt,a4paper]{article}
\usepackage[utf8]{inputenc}
\usepackage[margin=2.5cm]{geometry}
\usepackage{enumitem}
\usepackage{hyperref}
\usepackage{booktabs}
\usepackage{tabularx}
\usepackage{xcolor}

\definecolor{done}{RGB}{34,139,34}
\definecolor{todo}{RGB}{200,50,50}
\definecolor{question}{RGB}{0,90,180}

\newcommand{\done}[1]{\textcolor{done}{\textbf{DONE:}} #1}
\newcommand{\todo}[1]{\textcolor{todo}{\textbf{TODO:}} #1}
\newcommand{\question}[1]{\textcolor{question}{\textbf{Q:}} #1}

\title{Supervisor Meeting Agenda\\
\large 16 March 2026 --- 30 minutes}
\author{Antonio Galdes}
\date{}

\begin{document}
\maketitle

% ------------------------------------------------------------------
\section*{1. Progress Update (5 min)}
% ------------------------------------------------------------------

\begin{itemize}[leftmargin=*]
  \item \done{Tasks 1--8 code-complete. 163 unit tests passing, \texttt{make verify}
    green. Full GPU stack verified (RTX 4070 Ti Super): CARLA 0.9.16, ros2-bridge
    (Jazzy), training container (Humble/NGC PyTorch).}

  \item \done{Localisation architecture decided and fully implemented: Cartographer
    pure localisation against a frozen \texttt{.pbstream} map built from perimeter
    traffic cones. Empty lot = all cones visible = low EKF covariance (easy). Dynamic
    objects occluding cones = degraded scan match = high covariance (hard). Signal
    is correctly correlated with driving difficulty.}

  \item \done{One-time SLAM mapping workflow operational end-to-end:
    \texttt{make docker-mapping-drive LAYOUT=rectangle} drives the patrol-waypoint loop
    autonomously (2 loops, $\sim$90 s), Cartographer builds the map in real time,
    and \texttt{make docker-save-map LAYOUT=rectangle} serialises the
    \texttt{.pbstream} via the ROS 2 \texttt{/write\_state} service.
    \texttt{rectangle.pbstream} successfully generated today.}

  \item \done{NPC behaviour improved: patrol vehicles use a forward-cone proximity
    check (no longer stopping for pedestrians behind them), throttle raised to 0.6,
    pedestrians walk away from both ego and patrol NPCs using a combined repulsion
    vector at 1.0\,m/s escape speed.}

  \item \done{Stale \texttt{scripts/generate\_cartographer\_maps.py} (generated
    unusable \texttt{.pgm} files) removed. Replaced by the Cartographer-native
    \texttt{.pbstream} workflow above.}

  \item \todo{Mapping runs for \texttt{trapezoid} and \texttt{irregular\_a} still
    to be executed (same workflow, $\sim$90 s each).}

  \item \todo{First training run (\texttt{make docker-train}) not yet started ---
    mapping must complete for all three floor plans first.}
\end{itemize}

% ------------------------------------------------------------------
\section*{2. Discussion Points (15 min)}
% ------------------------------------------------------------------

\subsection*{2.1 Simulation Scenario Sufficiency}

The current simulation covers three floor plan geometries (rectangle, trapezoid,
irregular\_a), two weather presets (ClearNoon, HardRainNoon), NPC patrol vehicles,
random-walk pedestrians, and variable bay occupancy (0--100\%). Uncertainty varies
naturally: an empty lot with no occlusion produces low EKF covariance, while a
crowded lot with dynamic objects blocking cone returns produces high covariance.

\question{Is the current set of conditions (3 layouts $\times$ 2 weather presets
$\times$ variable occupancy/traffic/pedestrians) sufficient to produce a
statistically meaningful degradation curve for the evaluation chapter, or should
additional conditions (e.g.\ a third weather preset, additional floor plan variants)
be added before beginning training?}

\question{The covariance signal is driven entirely by cone occlusion from NPC cars
and pedestrians. Does this mechanism produce a sufficiently wide and smooth range
of uncertainty values to demonstrate the thesis claim (agent behaves more
conservatively at high uncertainty), or is an additional artificial noise injection
needed as a fallback?}

\bigskip
\subsection*{2.2 Cartographer Pure Localisation vs.\ SLAM During Training}

Cartographer now runs in pure localisation mode during training: the pose graph is
frozen, scans are matched against the pre-built cone map only, and no new submaps
are created. This is the architecture that produces the correct (positively
correlated) covariance signal.

\question{Is the one-time SLAM mapping step per floor plan (run once, reuse
\texttt{.pbstream} for all training episodes) academically acceptable, or does the
examiner need to be satisfied that the mapping step is fully automated and
reproducible? (It is: \texttt{make docker-mapping-drive} + \texttt{make docker-save-map}
are the only commands required.)}

\bigskip
\subsection*{2.3 Ablation Study Design}

The 2$\times$2 ablation compares: (1) vanilla PPO with no uncertainty in observation
or policy, (2) uncertainty in observation only, (3) uncertainty in policy output
only, and (4) the full method (both). All four share identical PPO hyperparameters,
sensor noise, and training conditions.

\question{Is a 2$\times$2 factorial ablation (4 conditions $\times$ 10 seeds each)
the right evaluation design for the contribution, or would a single full-method run
with a strong vanilla-PPO baseline be a simpler and equally convincing comparison?}

\question{Should the ablation include a condition where the Cartographer localisation
is replaced by ground-truth CARLA pose (zero uncertainty), to verify that the agent
actually uses the covariance signal rather than ignoring it?}

\bigskip
\subsection*{2.4 Reward Function}

The current reward combines a potential-based shaping term (distance to target bay),
a heading alignment bonus, a time penalty, a collision penalty, and a success bonus.
No reward tuning has been done yet --- this is Task 9, the first thing after the
initial training run.

\question{Any guidance on reward shaping priorities: should the agent first be
rewarded for reaching any bay successfully before introducing the uncertainty-aware
conservative behaviour objective?}

% ------------------------------------------------------------------
\section*{3. Timeline and Next Steps (5 min)}
% ------------------------------------------------------------------

\begin{table}[h]
\centering
\small
\begin{tabular}{@{}lll@{}}
\toprule
\textbf{Task} & \textbf{Est.\ Time} & \textbf{Status} \\
\midrule
Mapping runs (trapezoid, irregular\_a) & 0.5 days & Ready to run \\
First training run (\texttt{make docker-train}) & 1 day & Unblocked after mapping \\
Task 9: Reward function tuning & 1--2 weeks & After first run \\
Full ablation (4 baselines $\times$ 10 seeds) & 2--4 weeks & After Task 9 \\
Tasks 10--14: Evaluation + plots & 2 weeks & After training \\
Dissertation writing & Ongoing & Can start now \\
\bottomrule
\end{tabular}
\end{table}

\question{Does this timeline remain feasible given the September 2026 submission
deadline? Which tasks should be deprioritised if the schedule tightens?}

% ------------------------------------------------------------------
\section*{4. Dissertation Writing (5 min)}
% ------------------------------------------------------------------

\begin{itemize}[leftmargin=*]
  \item The methodology chapter can be drafted now --- the architecture is stable
    and all design decisions are resolved.
  \item The background chapter (EKF, evidential deep learning, PPO, Cartographer)
    is independent of experimental results and can be written in parallel with
    training.
  \item The experimental setup chapter is ready to write: floor plans, sensor suite,
    NPC configuration, and the Cartographer pure localisation workflow are all fixed.
\end{itemize}

\question{Any guidance on dissertation structure, expected length, chapter ordering,
or when to submit draft chapters for feedback?}

\question{Should the Cartographer mapping workflow (one-time SLAM session to build
the cone perimeter map) be described in the methodology chapter as part of the
system design, or relegated to an appendix as a deployment detail?}

\end{document}
